% TOP_PROBLEM: {"id": 3052, "title": "Roller Coaster permutations", "category_id": 2, "category": {"id": 2, "name": "combinatorics", "display_name": "Combinatorics", "description": "Counting problems, graph theory, discrete structures.", "slug": "combinatorics", "order_index": 2, "created_at": "2026-07-31T15:26:25.671Z"}, "status": "open"}
% FIRST_POSTED: null
% ENTRY_KIND: statement_only
% Imported from: batch01.tex; UnsolvedMath ID 3052
\documentclass[11pt]{book}
\usepackage{fontspec}
\setmainfont{Latin Modern Roman}
\setsansfont{Latin Modern Sans}
\usepackage{amsmath,amssymb,mathtools}
\usepackage{unicode-math}
\setmathfont{Latin Modern Math}
\usepackage{xurl}
\usepackage[hidelinks]{hyperref}
\newcommand{\sref}[2]{\hyperlink{source:#1:#2}{[S#2]}}
\newcommand{\cataloguescope}[1]{\paragraph{Mathematical question.}#1}
\begin{document}
% UnsolvedMath ID 3052; source catalogue code OPG-58213
\setcounter{section}{3051}
\section{Roller Coaster permutations}\label{3052}
{\small\sffamily UnsolvedMath ID 3052 · Combinatorics}
\subsection{Definitions and mathematical statement}

\paragraph{Shared notation.}$\mathbb N=\{1,2,\ldots\}$, and $\mathbb Z,\mathbb Q,\mathbb R,\mathbb C$ denote the integers, rationals, reals and complex numbers. Any other conventions are specified locally.


Let $S_n$ be the permutations of $[n]=\{1,\ldots,n\}$, written as sequences. An increasing (respectively decreasing) run is a maximal contiguous increasing (respectively decreasing) subsequence of length at least two; consecutive runs may share their turning-point entry. Write $i(\tau)$ and $d(\tau)$ for their numbers. If $\tau$ has length $m\geq2$, then
\[i(\tau)+d(\tau)=1+\sum_{j=2}^{m-1}\mathbf1\bigl\{(\tau_j-\tau_{j-1})(\tau_{j+1}-\tau_j)<0\bigr\}.\]
Let $X(\pi)$ consist of all subsequences selected by strictly increasing index sets of size at least three, and set
\[t(\pi)=\sum_{\tau\in X(\pi)}(i(\tau)+d(\tau)),\qquad\operatorname{RC}(n)=\{\pi\in S_n:t(\pi)=\max_{\sigma\in S_n}t(\sigma)\}.\]
The run convention is the one illustrated in the cited original research definition. \sref{3052}{3}
\paragraph{Statement recorded in the supplied source.}
For every $n\geq3$, are the following two conjectures true? First, if $n=2k$ then $|\operatorname{RC}(n)|=4$, while if $n=2k+1$ then $|\operatorname{RC}(n)|=2^j$ for some integer $0\leq j\leq k+1$. Second, for every $\pi\in\operatorname{RC}(n)$ and every $1\leq a\leq k$, is $\pi_a+\pi_{n-a+1}$ odd when $n=2k+1$, and equal to $2k+1$ when $n=2k$? \sref{3052}{2}
\subsection{Short English statement}
For permutations maximizing the total monotone-run count over all subsequences, determine whether their number is four in even length and a bounded power of two in odd length, and whether opposite-position entries satisfy the stated sum laws.
\subsection{Sources}
\begin{description}
\item[\hypertarget{source:3052:1}{S1}] ULAM.AI and the UnsolvedMath Contributors, \emph{UnsolvedMath: A Curated Collection of Open Mathematics Problems}, record 3052 (OPG-58213). CC BY 4.0. \url{https://huggingface.co/datasets/ulamai/UnsolvedMath}.
\item[\hypertarget{source:3052:2}{S2}] Ahmed, Tanbir; Snevily, Hunter S., Roller Coaster permutations, Open Problem Garden, node 58213. Complete problem, discussion and bibliography (including mathematical image alt text). \url{https://www.openproblemgarden.org/op/roller_coaster_permutations}.
\item[\hypertarget{source:3052:3}{S3}] William L. Adamczak, \emph{A Note on the Structure of Roller Coaster Permutations}, arXiv:1605.02415 (2016). Definition1.1 and Example1.2, page2; reference[1] identifies Ahmed--Snevily, BulletinICA68(2013),55--60. \url{https://arxiv.org/pdf/1605.02415}.
\end{description}
\label{end:3052}
\end{document}
